% ==============================================================================================================
% ================================================ INPUT =======================================================
% ==============================================================================================================


% ================================================ COVER =======================================================

% Logos located at Resources/logos/...
\def\coversubjectlogo{AST-h-black.png} 
\def\coversubjectlogoscale{0.6} 
%% [UC-TeX][INPUT] \specifictitlename se usa sola en Pages/Doc/short_cover.tex;
%% \largetitlename es el título completo de la portada.
\def\specifictitlename{Análisis de Espectros Estelares}
\def\largetitlename{Laboratorio 1: \specifictitlename}

% Select wether you want to include ORCID iDs and links
\newif\iforcidbool
% \orcidboolfalse   % or 
\orcidbooltrue

% Select citation and reference style:
%   \usenumcitetrue  -> numbered style [1], [2], ...
%   \usenumcitefalse -> default AASTeX author-year style (Author et al. YYYY)
%
% NOTE: after switching modes, delete .aux and .bbl before recompiling,
%       then run: pdflatex -> bibtex -> pdflatex -> pdflatex
\newif\ifusenumcite
% \usenumcitetrue   % or
\usenumcitefalse

%% [UC-TeX][INPUT] Referencias a una columna (ancho completo) o a dos columnas (por defecto).
% \onecolumnrefstrue   % o
\onecolumnrefsfalse

%% [UC-TeX][INPUT] Sintaxis de cada entrada: {ORCID}{Nombre completo}{Correo}{Afiliación}
%% Cualquier campo puede quedar vacío ({}). Correo y afiliación aparecen como nota al pie
%% en la portada. Si una lista queda vacía, su encabezado no se muestra.

% --- Authors ---
\def\coverauthors{
    {0009-0005-4685-695X}{Juan José Sánchez Medina}{}{},
    {}{Henrietta Leavitt}{henrietta.leavitt@estudiante.uc.cl}{Pontificia Universidad Católica de Chile},
    {}{Stephen Hawking}{}{}
}

% --- Professors ---
\def\coverprofessor{
    {0000-0000-0000-0000}{Isaac Newton}{}{}
}

% --- Assistants ---
\def\coverassist{
    {}{Marie Curie}{}{}
}

%% [UC-TeX][INPUT] Deje \docinitdate vacío ({}) para ocultar la fecha del laboratorio (p. ej. tareas).
% --- Lab Date ---
\def\docinitdate{6 de noviembre}
% --- Final Date ---
\def\docfinalmonth{noviembre}

% --- Subject Details ---
\def\subjectcode{AST067}
\def\subjecttext{Laboratorio de Astronomía}


% ================================================ HEADER ======================================================
\shorttitle{Informe Laboratorio 01}
\shortauthors{Grupo 2070}

%% Include dates for submitted, revised, and accepted.
%%\received{February 1, 2025}
%%\revised{March 1, 2025}
%%\accepted{\today}
%%
%% Indicate AAS Journal the manuscript was submitted to.
%%\submitjournal{PSJ}
%% Note that this command adds "Submitted to " the argument.
%%
%% You can add a light gray and diagonal water-mark to the first page 
%% with this command:
%% \watermark{text}
%% where "text", e.g. DRAFT, is the text to appear.  If the text is 
%% long you can control the water-mark size with:
%% \setwatermarkfontsize{dimension}
%% where dimension is any recognized LaTeX dimension, e.g. pt, in, etc.
%%%%%%%%%%%%%%%%%%%%%%%%%%%%%%%%%%%%%%%%%%%%%%%%%%%%%%%%%%%%%%%%%%%%%%%%%%%%%%%%
%%
%% Use this command to indicate a subdirectory where figures are located.
%%\graphicspath{{./}{figures/}}

% =============================================== PREAMBLE =======================================================

% --- Custom Imports ---
\usepackage{amsmath}
\usepackage{cancel}
\usepackage{lipsum}
\usepackage{multirow}
\usepackage{soul}
\usepackage[colorinlistoftodos, spanish]{todonotes}
\usepackage{siunitx}
%% [UC-TeX] Paquetes adicionales
\usepackage{amssymb}
\usepackage{float}       % [H] para figuras/tablas
\usepackage{placeins}    % \FloatBarrier
\usepackage{mhchem}      % \ce{^{60}Co}
\usepackage{tikz}
\usetikzlibrary{positioning, arrows.meta, calc}
\usepackage{physics}     % \dv, \pdv, \abs, ...
\AtBeginDocument{\RenewCommandCopy\qty\SI} % physics redefine \qty; se restaura el de siunitx

%% [UC-TeX] Estilos de código (listings): use [style=python] o [style=adql].
%% Comente las líneas de los lenguajes que no necesite.
%% [UC-TeX] Configuración base de listings (bloques de código).
%% Se carga desde Input.tex antes de los estilos de cada lenguaje (python.tex, adql.tex).
%% Uso:  \begin{lstlisting}[style=python, caption={...}] ... \end{lstlisting}
%%       \lstinputlisting[style=adql]{ruta/al/archivo.sql}

\usepackage{listings}
\usepackage{xcolor}

% --- Colores comunes ---
\definecolor{codegreen}{rgb}{0,0.6,0}
\definecolor{codegray}{rgb}{0.5,0.5,0.5}
\definecolor{codepurple}{rgb}{0.58,0,0.82}
\definecolor{backcolour}{rgb}{0.95,0.95,0.92}

% --- Nombres según idioma (es / en) ---
\makeatletter
\@ifclasswith{TeXAST}{es}{%
  \renewcommand{\lstlistingname}{Código}%
  \renewcommand{\lstlistlistingname}{Índice de código}%
}{%
  \renewcommand{\lstlistingname}{Listing}%
  \renewcommand{\lstlistlistingname}{List of Listings}%
}
\makeatother
% \lstlistoflistings   % índice de bloques de código (dentro del documento)

%% [UC-TeX] Estilo de código Python para listings. Requiere Config/Listings/base.tex.
%% Uso: \begin{lstlisting}[style=python] ... \end{lstlisting}

\lstdefinestyle{python}{
    language=Python,
    basicstyle=\ttfamily\small,
    keywordstyle=\color{blue},
    commentstyle=\color{green!50!black},
    stringstyle=\color{red},
    showstringspaces=false,
    breaklines=true,
    numbers=left,
    numberstyle=\tiny\color{gray},
    frame=single
}

%% [UC-TeX] Estilo de código ADQL/SQL (consultas a Gaia, SDSS, DESI, etc.) para listings.
%% Requiere Config/Listings/base.tex.
%% Uso: \begin{lstlisting}[style=adql] ... \end{lstlisting}

\lstdefinestyle{adql}{
    language=SQL,
    backgroundcolor=\color{backcolour},
    commentstyle=\color{codegreen},
    keywordstyle=\color{blue},
    numberstyle=\tiny\color{codegray},
    stringstyle=\color{codepurple},
    basicstyle=\ttfamily\small,
    breakatwhitespace=false,
    breaklines=true,
    captionpos=b,
    keepspaces=true,
    numbers=left,
    numbersep=5.5pt,
    showspaces=false,
    showstringspaces=false,
    showtabs=false,
    tabsize=2,
    literate={"}{{\char34}}1,
    morekeywords={CONTAINS, POINT, CIRCLE, POLYGON, REGION, BOX, TOP},
    deletekeywords={DEC}  % "dec" es una columna (declinación), no DECIMAL
}


%% [UC-TeX] Macros de astronomía (\ra, \dec, unidades siunitx, nombres de revistas).
%% [UC-TeX] Macros de astronomía. Se carga desde Input.tex (después de siunitx).

% --- Coordenadas ---
% Ascensión recta: \ra{hh}{mm}{ss}  ->  05^h 35^m 17^s
\newcommand{\ra}[3]{%
  \num[minimum-integer-digits=2]{#1}^{\text{h}}\, %
  \num[minimum-integer-digits=2]{#2}^{\text{m}}\, %
  \num[minimum-integer-digits=2]{#3}^{\text{s}}%
}
% Declinación: \dec{dd}{mm}{ss}  ->  -05° 23' 28''
\newcommand{\dec}[3]{%
  \num{#1}^{\circ}\, %
  \num[minimum-integer-digits=2]{#2}'\, %
  \num[minimum-integer-digits=2]{#3}''%
}

% --- Unidades siunitx: \qty{10}{\parsec}, \unit{\hubblec} ---
\DeclareSIUnit\angstrom{\textup{\AA}}
\DeclareSIUnit\parsec{\textup{pc}}
\DeclareSIUnit\year{\textup{yr}}
\DeclareSIUnit\lightyr{\textup{ly}}
\DeclareSIUnit\jansky{\textup{Jy}}
\DeclareSIUnit\pixel{\textup{px}}
\DeclareSIUnit\mas{\textup{mas}}
\DeclareSIUnit\mag{\textup{mag}}
\DeclareSIUnit\maggie{\textup{mgy}}
\DeclareSIUnit\hubblec{\kilo\meter\per\second\per\mega\parsec}

% --- Varios ---
\newcommand{\cfeh}{$\left[ \rm Fe / H \right]$}
\newcommand{\tick}{$\checkmark$}

% --- Nombres completos de revistas en la bibliografía (en vez de ApJ, ApJL, Nature) ---
% Comente estas líneas para volver a las abreviaturas de AASTeX.
\def\apjl{The Astrophysical Journal Letters}
\def\apj{The Astrophysical Journal}
\def\nat{Nature}


% --- Custom Commands ---
\newcommand{\vdag}{(v)^\dagger}
\newcommand\aastex{AAS\TeX}
\newcommand\latex{La\TeX}
\newcommand{\diff}{\mathop{}\!\mathrm{d}}
\newcommand{\cint}{%
  \mathop{%
    \vcenter{\hbox{\(\displaystyle\int\mkern-18mu C\)}}%
  }\nolimits%
}

%% [UC-TeX] Encabezado centrado dentro de una deluxetable (envuelve \cutinhead).
\newcommand{\headgroup}[1]{\cutinhead{\makebox[0pt]{#1}}}

\newcommand{\todoerror}[2][]{\todo[color=red!70, inline, #1]{#2}}
\newcommand{\todowarn}[2][]{\todo[color=orange, inline, #1]{#2}}
\newcommand{\todoremark}[2][]{\todo[color=yellow!80!black, inline, #1]{#2}}
\newcommand{\todoshort}[2][]{\todo[color=yellow!80!black, #1]{#2}} %% [UC-TeX] nota al margen
\newcommand{\todohint}[2][]{\todo[color=green!20, inline, #1]{#2}}
\newcommand{\tododone}[2][]{\todo[color=blue!50, inline, #1]{#2}}

\makeatletter
\let\frontmatter@title@above=\relax
\makeatother